\documentclass[10pt,a4paper]{amsart}
\usepackage[margin=19mm]{geometry}
\usepackage{amsmath,amssymb,mathtools}
\usepackage[scr=rsfs,cal=euler]{mathalfa}
\usepackage{xfrac}
\usepackage[unicode,hidelinks,bookmarks=false]{hyperref}
\newcommand{\ie}{i.e.\ }
\newcommand{\cf}{cf.\ }
\newcommand{\resp}{respectively\ }
\newcommand{\Subsection}{\subsection}
\providecommand{\nobreakdash}{\nobreak\hbox{-}}
\setlength{\emergencystretch}{2em}
\multlinegap0pt
\usepackage{bm}
\usepackage{bbm}
\usepackage{dsfont}
\usepackage{xspace}
\usepackage{cancel}
\usepackage{mathtools}
\usepackage{amsthm}
\makeatletter
\@ifundefined{assumption}{\newtheorem{assumption}{Assumption}}{}
\@ifundefined{theorem}{\newtheorem{theorem}{Theorem}}{}
\makeatother
\title[Synchronization-Free Algebraic Fingerprints for Large Langua]{Synchronization-Free Algebraic Fingerprints for Large Language Models: From Autoregressive to Diffusion Models}
\author{Jaroslaw Janas, Josef Pieprzyk, Pawel Morawiecki}
\date{Source: arXiv:2607.16648v1 (CC BY 4.0)}
\begin{document}
\maketitle

\noindent\emph{Source and licence.} Synchronization-Free Algebraic Fingerprints for Large Language Models: From Autoregressive to Diffusion Models. Version arXiv:2607.16648v1 (CC BY 4.0), distributed under the Creative Commons Attribution 4.0 International licence, as recorded in the arXiv licence metadata for this exact version and shown on the abstract page https://arxiv.org/abs/2607.16648v1. Full rights notice: licenses/arxiv-2607.16648-CC-BY-4.0.txt.\par\smallskip

\noindent\textit{Editorial scope.} Counted unit: the Theorem whose source label is \texttt{None} in the article (source0.tex lines 2164--2182, source label \texttt{None}), delivered together with the article's own complete original proof of it (source0.tex lines 2184--2230, one \texttt{proof} environment of 47 lines). No mathematics is written, repaired or re-derived: statement and proof are the article's own text line for line. Required context retained uncounted: ass:probability (lines 1630--1665). Every definition, display and label of those lines is retained unchanged. Omitted from this package and not redistributed: 1--1629, 1666--2163, 2183--2183, 2231--2472. Nothing inside a retained excerpt is omitted or altered: every retained body is delivered exactly as the article's lines are written, so no omission receipt of any kind accompanies this package. Citations: the retained excerpts carry no citation command at all, so no bibliography entry of the article is transcribed and no reference is redistributed. Reference adaptation: none was needed and none was made; every reference inside the delivered unit resolves inside it, so no unresolved reference or citation is produced. Printed slips kept literal: no printed word, symbol, hypothesis or number of the counted statement or of its proof was altered. Layout and class: the article's own class is \texttt{llncs} with packages including float, graphicx, fontenc, url, mathtools, nicefrac, tcolorbox, braket, booktabs, xcolor, framed, color, amsmath, amssymb; the wrapper is the lane's pinned \texttt{amsart} editorial wrapper at 10pt on A4, which restates the theorem-like environments the retained text needs and adds the article's own macro definitions that the retained text needs (none), with extra wrapper packages (bm, bbm, dsfont, xspace, cancel, mathtools). The delivered numbering is the unit's own. Assets: no retained span contains a figure environment, a graphics-inclusion command, a PDF-page inclusion, a table or a TikZ picture; no image file is copied into this package, the package declares no asset dependency and no image file is redistributed with it. Licence: version arXiv:2607.16648v1 of this article is distributed under the Creative Commons Attribution 4.0 International licence, as recorded in the arXiv licence metadata for this exact version and shown on the abstract page https://arxiv.org/abs/2607.16648v1. A full rights notice is delivered at licenses/arxiv-2607.16648-CC-BY-4.0.txt. Characters: every retained excerpt is pure ASCII, so the delivered engine, XeLaTeX with its default Latin Modern text font, reports no missing character.
\begin{assumption}
\label{ass:probability}
The colour selected for each generated token is treated as an
independent random variable.
The probability that the \emph{left constraint} is satisfied is denoted
by $\varepsilon$, where
$
\frac12 \le \varepsilon \le 1.
$
The parameter $\varepsilon$ models the effectiveness of the watermark
embedding mechanism. Modern LLM watermarking algorithms bias the logits
towards the desired colour class before sampling the next token.
Consequently, whenever both colour classes contain sufficiently many
candidate tokens, the probability of selecting a token from the desired
colour class can be made arbitrarily close to one. On the other hand,
there are rare situations in which the language model offers no suitable
token of the required colour, making successful embedding impossible.
For this reason, $\varepsilon$ should be regarded as an average success
probability over all generated tokens.
Throughout the paper we analyse the algorithms for an arbitrary
$\varepsilon$, while the conservative choice
$$
\varepsilon=\frac12
$$
corresponds to unbiased sampling and therefore represents the worst
practical case.
The \emph{right constraint} depends only on the cryptographic hash
function. Under the standard assumption that the hash behaves as a
pseudorandom function, its output is uniformly distributed, giving
$$
\Pr[\text{right constraint holds}]
=\frac12.
$$
Furthermore, the left and right constraints are assumed to be
independent.
\end{assumption}

\begin{theorem}
Under Assumption~\ref{ass:probability}, every active boundary advances by one position with
probability
$
\frac12,
$
during each iteration. Consequently, the expected displacement after
$t$ iterations equals
$
\frac{t}{2},
$
while the variance equals
$
\frac{t}{4}.
$
Thus the commitment boundary performs a biased random walk whose drift
is independent of the watermark embedding probability
$\varepsilon$.
\end{theorem}

\begin{proof}
A boundary movement occurs precisely when exactly one of the two
constraints holds.
Since
$
\Pr[L\wedge\overline{R}]
=
\frac{\varepsilon}{2}
$
and
$
\Pr[\overline{L}\wedge R]
=
\frac{1-\varepsilon}{2},
$
their sum equals
$
\frac{\varepsilon}{2}
+
\frac{1-\varepsilon}{2}
=
\frac12.
$
Each successful movement shifts the boundary by one token.
The number of successful movements after $t$ iterations therefore
follows the binomial distribution
$
X_t\sim\mathrm{Binomial}
\left(
t,\frac12
\right).
$
Hence
$
E[X_t]
=
\frac{t}{2},
$
and
$
\mathrm{Var}(X_t)
=
\frac{t}{4},
$
which proves the theorem.
\hfill$\Box$
\end{proof}

\end{document}
